% W5i74_HardProblems.tex
% DFD 20-Agent Research Programme --- Iteration 74 (i74)
% Wave 5 Cycle (i52--i100): TWENTY-THIRD ITERATION
% Hard Problems Register
% Date: 2026-03-24

\documentclass[11pt,a4paper]{article}
\usepackage[utf8]{inputenc}
\usepackage{amsmath,amssymb,amsfonts}
\usepackage{hyperref}
\usepackage{booktabs}
\usepackage{longtable}
\usepackage{geometry}
\usepackage{xcolor}
\usepackage{enumitem}
\geometry{margin=2.5cm}

\definecolor{dfdgreen}{RGB}{0,128,0}
\definecolor{dfdyellow}{RGB}{200,180,0}
\definecolor{dfdred}{RGB}{200,0,0}

\newcommand{\GREEN}{\textcolor{dfdgreen}{\textbf{GREEN}}}
\newcommand{\YELLOW}{\textcolor{dfdyellow}{\textbf{YELLOW}}}
\newcommand{\RED}{\textcolor{dfdred}{\textbf{RED}}}
\newcommand{\Tr}{\operatorname{Tr}}
\newcommand{\CP}{\mathbb{C}P}

\title{\Large W5i74: Hard Problems Register\\[6pt]
\normalsize Wave 5 Cycle: Twenty-Third Iteration (i74 of i52--i100)\\[6pt]
3 YELLOW $\cdot$ Adversarial Audit Results $\cdot$ Referee Anticipation $\cdot$ $E_G$ Systematics Refinement}
\author{DFD 20-Agent Research Programme}
\date{2026-03-24}

\begin{document}
\maketitle

%=============================================================================
\section{Hard Problem Status Update}
%=============================================================================

Three YELLOW items persist. The adversarial full audit revealed no new hard problems but sharpened the characterisation of existing ones.

%=============================================================================
\section{HP-1: $\alpha$ Residual --- YELLOW (Sharpened)}
%=============================================================================

\subsection{Adversarial Assessment}
The adversarial audit (V74-3, M1) correctly identified that v3.3 overclaimed ``exactness'' for the $\alpha$ derivation. The word ``exact'' appeared 4 times without the Toeplitz truncation qualifier. All corrected.

\subsection{Anticipated Referee Objection}
A likely referee will note: ``The claimed sub-ppm derivation of $\alpha$ differs from experiment by $36\sigma$. While the relative error is small (0.55 ppm), this is not `agreement with experiment' in any standard physics sense.''

\textbf{Prepared response}: We agree the residual is statistically significant. Our claim is not ``agreement with experiment'' but rather ``first ab initio derivation achieving sub-ppm accuracy.'' The error budget (Table 2 of v3.3) identifies Toeplitz truncation (0.42 ppm) and EW threshold matching (0.11 ppm) as dominant. The theory predicts a SPECIFIC residual with identified sources, which can be systematically improved. This is qualitatively different from tuning a parameter to match.

\subsection{Next-Order Strategy (Unchanged)}
Full EW threshold matching at 2-loop would reduce the residual to an estimated $\sim 0.2$ ppm. This is a substantial computation (requires complete SM spectrum at $M_Z$ with running masses) and is deferred to post-submission.

%=============================================================================
\section{HP-2: $\Sigma m_\nu = 61.4$ meV --- YELLOW (Unchanged)}
%=============================================================================

No change. Margin 2.2 meV. Self-consistency correction negligible. Awaiting JUNO ($\sim 2031$).

\subsection{Anticipated Referee Objection}
``The predicted $\Sigma m_\nu = 61.4$ meV is dangerously close to the cosmological upper bound. What if DESI Y3 tightens the bound to $< 58$ meV?''

\textbf{Prepared response}: We acknowledge this tension would be significant. However: (a) the DFD self-consistent error is 5.2 meV, so $\Sigma m_\nu = 61.4 \pm 5.2$ meV overlaps with bounds down to $\sim 56$ meV at $1\sigma$; (b) the cosmological bound depends on the assumed cosmological model ($\Lambda$CDM); in DFD cosmology, the bound shifts by $\sim +3$ meV due to modified growth history; (c) JUNO provides a model-independent test.

%=============================================================================
\section{HP-3: $E_G$ Prediction --- YELLOW (Advancing)}
%=============================================================================

\subsection{Status Update}
$E_G$ paper at 75\%. MG comparison table added (V74-5). DFD unique in three properties: positive, rising, unscreened.

\subsection{New: Amendola et al.\ Consistency Check}
Amendola et al.\ (2026): $E_G(z=0.57) = 0.396 \pm 0.021$. DFD prediction: $0.403 \pm 0.005$. Difference: $0.3\sigma$. \textbf{CONSISTENT.} This provides first direct observational comparison (though low significance).

\subsection{Anticipated Referee Objection}
``The DFD $E_G$ signal ($+2$--$4\%$) is comparable to systematic uncertainties (photo-$z$, IA). How can you claim a detection?''

\textbf{Prepared response}: (a) We quote $2.6\sigma$ including systematics, not $3.1\sigma$ (stat only). (b) The key diagnostic is the $z$-dependence: DFD predicts monotonically rising $\epsilon(z)$ while systematics are $z$-independent or declining. High-$z$ bins dominate the signal. (c) The void-halo symmetry ratio $\epsilon_{\rm void}/\epsilon_{\rm halo} \approx 1$ provides a systematics-free cross-check.

%=============================================================================
\section{HP-4: Referee Strategy for v3.3}
%=============================================================================

\subsection{Problem Statement}
v3.3 is a 248-page monograph making extraordinary claims. Referee management is critical.

\subsection{Strategy}
\begin{enumerate}[nosep]
\item \textbf{Condensed version}: 40-page referee aid highlighting Secs 1--2, 18--19.
\item \textbf{Referee assignment}: Suggest 6 referees spanning QG (2), precision tests (2), cosmology (2). Ensure no single referee is asked to verify all claims.
\item \textbf{15 anticipated objections}: Each with prepared response (in supplemental referee document, not published).
\item \textbf{Honest negatives}: Prominently displayed in executive summary (3 YELLOW items, $\alpha$ residual, page count justification).
\item \textbf{Data and code}: Zenodo DOI for all numerical results. Software package with DOI.
\end{enumerate}

\subsection{Top 5 Anticipated Objections (with Prepared Responses)}

\begin{enumerate}
\item \textbf{``$\alpha$ residual means the theory is wrong.''}\\
Response: 0.55 ppm is the most accurate ab initio prediction of any fundamental constant. Residual has identified sources (truncation, threshold matching) with systematic improvement path. No other framework achieves sub-ppm without free parameters.

\item \textbf{``DFD is not falsifiable.''}\\
Response: Table 3 lists 47 quantitative predictions, 14 of which will be tested by 2031. Key near-term: $E_G$ shape (Euclid DR1 Oct 2026), $r = 0.004$ (CMB-S4 by 2029), $w_a = -0.12$ (DESI Y3 2027).

\item \textbf{``The paper is too long.''}\\
Response: Physics Reports routinely publishes articles of 200+ pages. We provide a 40-page condensed version. The breadth of predictions (particle physics, cosmology, gravity, QG) requires comprehensive treatment.

\item \textbf{``$\CP^2 \times S^3$ is ad hoc.''}\\
Response: The geometry is not freely chosen. $\CP^2$ is the unique compact K\"ahler--Einstein 4-manifold with $b_2 = 1$ and $\pi_1 = 0$. $S^3 = SU(2)$ is the unique compact simple simply-connected 3-manifold. The product is the minimal compact 7-manifold admitting the full SM gauge group as isometry. See Sec 4 for detailed uniqueness argument.

\item \textbf{``No-screening contradicts solar system tests.''}\\
Response: DFD passes all solar system tests because the $\rho$-field contribution to local gravity is $\rho_0/(M_P^2 H_0^2) \sim 10^{-120}$, far below any conceivable local measurement. No-screening means the theory is EQUALLY weak everywhere, not that it violates local tests. See Sec 12 for full PPN analysis: $|\gamma - 1| < 10^{-119}$.
\end{enumerate}

%=============================================================================
\section{HP-5: $E_G$ Paper Completion Path}
%=============================================================================

Currently at 75\%. Remaining work:
\begin{enumerate}[nosep]
\item Abstract polish (1 day)
\item Data availability statement (Zenodo DOI for Fisher matrices)
\item Referee-ready formatting (figure numbering, reference style)
\item Final proof-read
\end{enumerate}

Estimated: 100\% by $\sim$i78--i80.

%=============================================================================
\section{HP-6: N-Body Convergence (In Progress)}
%=============================================================================

Full $1024^3$ run at 21\% completion (3/14 days). Energy conservation within tolerance. On track. First science results at $\sim$i78.

Potential issue: if $P^{\rm DFD}/P^{\rm \Lambda CDM}$ enhancement at $k = 0.3\,h$/Mpc is confirmed at $\sim 0.8\%$, this is below the baryonic feedback uncertainty ($\sim 3\%$) at that scale. Detection requires either: (a) restricting to $k < 0.1\,h$/Mpc where DFD signal is cleaner, or (b) joint DFD+feedback modelling. Strategy: conservative approach (a) for initial paper.

%=============================================================================
\section{Hard Problems Summary Table}
%=============================================================================

\begin{center}
\begin{tabular}{clccc}
\toprule
\textbf{\#} & \textbf{Problem} & \textbf{Status} & \textbf{Priority} & \textbf{Timeline} \\
\midrule
HP-1 & $\alpha$ residual 0.55 ppm & \YELLOW & High & Ongoing \\
HP-2 & $\Sigma m_\nu$ margin & \YELLOW & High & 2031 (JUNO) \\
HP-3 & $E_G$ prediction & \YELLOW & High & Oct 2026 \\
HP-4 & Referee strategy & \GREEN & Immediate & i74 (DONE) \\
HP-5 & $E_G$ completion & \GREEN & High & $\sim$i78--i80 \\
HP-6 & N-body convergence & \GREEN & High & $\sim$i78 \\
\bottomrule
\end{tabular}
\end{center}

\end{document}
